% This is a template for your annotated bibliography. The only
% things you need to edit in this file are the title and author. You'll
% write the annotations themselves in the "annote" fields of the entries in
% bibliography.bib.
%
% To compiling documents on the command line, you'll run the following:
% latexmk -pdf annotated_bibliography.tex
%
% If you there is an error in the .bib file, you might need to delete the
% intermediate .aux, .bbl, and .blg files that are generated during the
% compilation process.

\documentclass[12pt]{article}
\usepackage{setspace}
\singlespace
\usepackage[left=1in,right=1in,top=1in,bottom=1in]{geometry}

\title{\textbf{Annotated Bibliography}}
\author{Your Name}

\begin{document}

\maketitle

% Using this makes it so that all the references show up regardless of whether
% or not they were actually cited in the rest of the document. Since we are
% only displaying the bibliography we won't be citing anything anywhere, so we
% need this. 
% https://www.mdpi.com/2076-3417/13/6/3895

\cite{su13031551}
%Simulation, Optimization, and Machine Learning in Sustainable Transportation Systems: Models and Applications 
The article provides a summary of techniques presently used in sustainable development within the scientific community.
The article begins by discussing the importance of sustainable development, the frequency of various techniques used to 
analyze sustainable development, and important issues in sustainable transportation, such as logistics and 
environmental issues. It then dives into the techniques that are currently being used to analyze sustainable
development, which includes simulation, optimization, machine learning, and fuzzy set analysis. The Optimization section covered
the history of optimization used in transportation and the variety of applications for optimization, including
in the difference between urban transportation systems and in global/regional transportation systems. Applications in simulation
focus on the growing integration of simulation and optimization models and on usages in logistics, traffic congestion, and on
the usage of demand-based transportation systems (ride sharing applications). The machine learning section primarily summarizes
a paper that uses a random forest model to understand patterns in a bike sharing transportation system, as well as touches on a
variety of techniques (random forest, logistic regression, neural networks, support vector machines, decision trees,
heuristics, multi-agent systems, and multi-variate analysis) for a variety of analytical, traffic-related outcomes.


\cite{su13031551}
%GeoAI: Integration of Artificial Intelligence, Machine Learning, and Deep Learning with GIS


\cite{ijgi11020085}
%Emerging Technologies for Smart Cities’ Transportation: Geo-Information, Data Analytics and Machine Learning Approaches 


% Use of machine learning in understanding transport dynamics of land use and public transportation in a developing city
\nocite{*}

\bibliographystyle{plain-annote}
\bibliography{bibliography.bib}

\end{document}
